\subsection{Den tekniske gennemspilning}
Den tekniske del af Oberstprøven handler i høj grad om at optimere forestillingens flow. Derfor udføres alle sceneskift to gange for at sikre at de kan afvikles konsekvent og i tide. Her tester man også forestillingens \textbf{cues}.
\subsubsection{Cues}
Da halvdelen af revyens indslag er sange, har bandet en vigtig rolle i at give revyen et godt \textit{flow}, dvs. publikums følelse af at blive båret godt videre fra det ene nummer til det andet. Sange starter typisk med, at bandet modtager et cue, hvorpå bandet sætter sangen i gang. Med et \textit{cue} forstås enhver form for signal, der informerer bandet  (eller en anden aktør) om at de skal gøre noget.
\\
\textbf{Bandets vigtigste tekniske opgave er, at have styr på bandets cues.}
\\
Derfor skal bandet for hver sang der spilles have følgende afklaret:
\begin{enumerate}
    \item Hvad er bandets cue, og hvem giver det?
    \item Hvilket bandmedlem starter sangen?
    \item Hvordan blive bandets cue kommunikeret til den som starter sangen?
\end{enumerate}
Typisk er svaret på disse spørgsmål:
\begin{enumerate}
    \item Bandets cue er, at sangeren rækker en hånd i vejret fra scenen.
    \item Trommeslageren starter sangen ved at tælle ind
    \item Et bandmedlem der både kan se scenen og trommeslageren giver signal til trommeslageren når cue'et modtages.
\end{enumerate}

\textbf{Konkret vil jeg foreslå,} at der er éen cue-ansvarlig for hele forestillingen. Dette minimerer forvirring for alle parter, og reducerer mental load. Den cue-ansvarlige burde notere cues i sin aktoversigt. Jeg vil også anbefale at den typiske løsning indikeret ovenfor (cue fra scene, cue kommunikeres til trommeslager, trommeslager tæller ind) som default, med mulighed for at justere efter kunstneriske og tekniske behov.

\textbf{Den optimale position for den cue-ansvarlige} er lige ved stigen til bandscenen. Herfra kan man se hele scenen, trekanten (bag scenetæppet, stage left), TeXnikken og det mest af bandet. På den måde kan man modtage og videregive nærmest alle cues. I særdeleshed er man tæt på trommeslageren, hvilket gør videregivelse af cues nemmere. Under Oberstprøven er det også ekstremt nyttigt at kunne kommunikere verbalt med trommeslageren, hvilket nærmest kun er muligt fra denne position. Dette kan også være en stor fordel hvis noget går galt under en forestilling.

\subsubsubsection{Cues der fungerer godt}
\textbf{Person rækker en arm i verjet.} Dette er tydeligt og utvetydigt, og varer ved så længe det er nødvendigt. Således er det robust overfor, at bandets cue-ansvarlige kigger væk et øjeblik.\\
\textbf{Person vender den røde lampe på mikrofon op imod bandet.} Den røde LED-lampe på de håndholdte mikrofoner er tydelig i mørket, og kan kun ses når den vendes rigtigt, og scenegænger ikke holder for den. Samtidig kan den ikke ses af publikum, hvis den vender imod bandet.\\
\textbf{Hvis sangen starter med solo musiker, kan denne tage cue direkte.} Hvis sangen f.eks. startes med klaverakkompagnement alene er der ingen grund til at sende cue til trommeslager. Lad da pianisten starte selv.

\subsubsubsection{Cues der ofte foreslås, men som fungerer dårligt}
\textbf{Når person sætter sig ned på scenen.} Efterspørges typisk hvis sangen starter med dansere (og evt. sanger) i en siddende/liggende pose. Problemet her er, at cue'et can overses, og derefter ikke gives igen. Et visuelt cue burde være stabilt indtil det modtages.\\
\textbf{Person bag bagtæppet giver signal.} Når der er mørkt i salen kan man overhovedet ikke se noget der sker bag bagtæppet fra bandscenen, ej heller hvis personen der giver tegn, står lige ved stigen.


\subsubsection{Lyd}
Sørg for, at der er hul igennem på alle instrumenter der bruges på et givet nummer. Vær særligt opmærksom på instrumenter der kun bruges på nogle numre, da disses kanaler kan være muted som udgangspunkt, og derfor manuelt skal unmutes på de relevante numre.\\
Det kan være gavnligt at tage sit monitorsystem af, og høre salens lyd med jævne mellemrum. Dette kan potentielt afdække tekniske problemer med lyd.\\
(TeXniken har på pulten gemt en scene pr. nummer i revyen. Typisk vil der være lavet en standard-indstilling som bruges som default på alle numre. Alligevel skal der laves et separat mix af hver scene. Det betyder også at der skal åbnes for de relevante kanaler på hver scene, hvis de er mutet by default.)

\subsubsection{"Runden"}
Efter hvert sceneskift vil den ansvarlige for afvikling af Oberstprøven tage en "runde" hvor de forskellige afdelinger får mulighed for at italesætte problemer. I interesse af at spare tid, er det gavnligt hvis bandet kan skelne mellem ting der kan løses internt i bandet, og ting der skal tages i plenum. Spørgsmål som "hvor mange takter er introen" klares internt. Spørgsmål som "er der lyd i min synth" skal udtrykkes i plenum, så Lyd kan svare.

\subsubsection{Pro tips}
\begin{itemize}
    \item Sørg for masser af snacks og vand på bandscenen.
    \item Lav kopholdere på bandscenen ved at tape Caféen? eller RKG-kopper på rælingen eller vægen.
    \item Bed scenemester om at sørge for, at alle rekvisiter der stilles bag bagtæppet har en ansvarlig der fjerner denne i tilfælde af brand. Således kan bandet også komme ned af bandscenen hvis det bliver nødvendigt.
\end{itemize}